\chapter{The Central Limit Theorem}

In Chapter~20 we proved the Law of Large Numbers: averaging many independent
copies of a random quantity pulls the sample mean $\bar X_n$ toward $\mu$.
The gap $\bar X_n - \mu$ collapses to zero. But what does the leftover
fluctuation look like when magnified by $\sqrt{n}$?

It turns out the shape does not depend on the original distribution. Whether you
start with dice rolls, coin flips, or exponential waiting times, adding
independent copies and standardizing the scale yields the exact same smooth,
symmetric bell curve. Every time.

This is the Central Limit Theorem (CLT), and it is foundational to probability and statistics.
It explains why the normal distribution appears across applications: measurement error,
survey sampling, A/B testing, physical experiments, and confidence intervals in machine learning
all rely on the same mathematical mechanism. We will state it carefully, watch it happen in simulation,
understand \emph{why} it is true (with both the intuitive moment generating function
argument and the rigorous characteristic function proof), learn the practical tools
that make it usable (standardization, the continuity correction), and find out how fast
it takes effect (Berry--Esseen). Along the way, we will see why the bell curve is
geometrically natural: every sum of independent variables is governed by convolution,
which progressively smooths probability densities as more terms are added.

\chapterline

\section{The phenomenon, first}
\label{ch21:phenomenon}

Start not with an abstract theorem but with a concrete progression. Roll a single fair
six-sided die. The outcome is uniform on $\set{1,2,3,4,5,6}$ -- a flat, boxy
distribution, six bars of equal height. Nothing bell-shaped about it.

Now roll \emph{two} dice and look at the sum. The sum runs from $2$ to $12$, and its
distribution is no longer flat: it is a triangle, peaked at $7$, because there are more
ways to make a $7$ ($1{+}6, 2{+}5, \dots$) than to make a $2$ ($1{+}1$ only). Already
the act of summing has pulled probability toward the middle.

Roll three dice, four dice, ten dice: the triangle rounds into a smooth dome, the
sharp corners vanish, and when centered and scaled to hold its variance constant, the
distribution settles toward a single fixed shape -- the standard normal bell curve.
The die was just the starting point; the repeated addition did the work.

\begin{intuitionbox}
A single random draw carries all the quirks of its parent distribution -- its skewness,
discreteness, or tail behavior. But when you add many independent draws together, those
finer details wash out, and the sum is governed almost entirely by just two numbers: its
mean (where it is centered) and its variance (how wide it spreads). What remains is the
distribution determined by mean and variance alone: the normal distribution.
\end{intuitionbox}

Figure~\ref{ch21:fig-dice} shows the dice story numerically. Watch the boxy uniform
become a triangle become a bell.

\begin{figure}[H]
\centering
\begin{tikzpicture}
\begin{axis}[
  width=13cm, height=5.2cm,
  ybar, bar width=3pt,
  xlabel={value of the sum}, ylabel={probability},
  xmin=0, xmax=37, ymin=0,
  scaled y ticks=false,
  yticklabel style={/pgf/number format/fixed, /pgf/number format/precision=2},
  title style={font=\small},
  title={Sum of $n$ fair dice: $n=1$ (flat), $n=2$ (triangle), $n=5$ (bell)},
  legend style={font=\footnotesize, at={(0.98,0.97)}, anchor=north east},
]
% n=1: uniform on 1..6, prob 1/6 ~ 0.1667
\addplot+[draw=sourceblue, fill=sourceblue!40] coordinates {
 (1,0.1667)(2,0.1667)(3,0.1667)(4,0.1667)(5,0.1667)(6,0.1667)};
\addlegendentry{$n=1$}
% n=2: triangular, sum 2..12, counts/36
\addplot+[draw=quotientorange, fill=quotientorange!40] coordinates {
 (2,0.0278)(3,0.0556)(4,0.0833)(5,0.1111)(6,0.1389)(7,0.1667)
 (8,0.1389)(9,0.1111)(10,0.0833)(11,0.0556)(12,0.0278)};
\addlegendentry{$n=2$}
% n=5: sum 5..30, exact counts/7776, symmetric about mean 17.5
\addplot+[draw=densityred, fill=densityred!40] coordinates {
 (5,0.00013)(6,0.00064)(7,0.00193)(8,0.00450)(9,0.00900)(10,0.01620)
 (11,0.02636)(12,0.03922)(13,0.05401)(14,0.06944)(15,0.08372)
 (16,0.09452)(17,0.10031)(18,0.10031)(19,0.09452)(20,0.08372)
 (21,0.06944)(22,0.05401)(23,0.03922)(24,0.02636)(25,0.01620)
 (26,0.00900)(27,0.00450)(28,0.00193)(29,0.00064)(30,0.00013)};
\addlegendentry{$n=5$}
\end{axis}
\end{tikzpicture}
\caption{The distribution of the sum of $n$ fair dice. With $n=1$ it is uniform (flat).
With $n=2$ it is a perfect triangle. With $n=5$ it is already strikingly bell-shaped,
even though the underlying die has nothing bell-shaped about it.}
\label{ch21:fig-dice}
\end{figure}

The geometric engine driving this transformation is \emph{convolution}---the
operation that computes the distribution of a sum of independent variables by sliding
one density across the other and integrating their overlap:
\[
  (f * g)(z) = \int_{-\infty}^{\infty} f(x)\, g(z - x)\,\dd x.
\]
Think of convolution as a continuous moving average: as $z$ shifts, the inverted profile
$g(z-x)$ glides across $f(x)$, blending and smearing local variations. Each convolution
acts as an integrator, smoothing the result and raising its smoothness class by one level
(where $C^0$ means continuous with no breaks, $C^1$ means continuous slope with no sharp
corners, and $C^\infty$ means infinitely smooth with every derivative well-behaved):
\begin{itemize}
  \item At $n=1$, the uniform density is discontinuous ($C^{-1}$ with jump steps at the boundaries).
  \item At $n=2$, convolving two step functions yields a continuous piecewise-linear triangle ($C^0$).
  \item At $n=3$, convolving again yields a quadratic spline ($C^1$ with continuous slope)---a curve stitched together from parabolic arches with matching slopes, rounding off the sharp peak.
  \item By $n=4$ and beyond, the curve has continuous curvature ($C^2$) and rapidly approaches the infinite smoothness ($C^\infty$) of the Gaussian bell.
\end{itemize}
Figure~\ref{ch21:fig-convolution-smoothing} demonstrates this progressive smoothing of Irwin--Hall distributions (the classic name for sums of independent uniform random variables on $[0,1]$).

\begin{figure}[H]
\centering
\begin{tikzpicture}
\begin{axis}[
  name=p1,
  width=7.2cm, height=4.2cm,
  xlabel={$x$}, ylabel={density},
  xmin=-0.5, xmax=4.5, ymin=0, ymax=1.15,
  title style={font=\small},
  title={Step 1: Uniform $n=1$ ($C^{-1}$, jump steps)},
]
\addplot[sourceblue, very thick, fill=sourceblue!20] coordinates {
  (-0.4,0) (0,0) (0,1) (1,1) (1,0) (4.4,0)
};
\node[font=\footnotesize, sourceblue, anchor=south west] at (axis cs:0.1,0.5) {flat block};
\end{axis}

\begin{axis}[
  name=p2, at={($(p1.east)+(1.2cm,0)$)}, anchor=west,
  width=7.2cm, height=4.2cm,
  xlabel={$x$}, ylabel={density},
  xmin=-0.5, xmax=4.5, ymin=0, ymax=1.15,
  title style={font=\small},
  title={Step 2: Sum of $2$ ($C^0$, continuous triangle)},
]
\addplot[quotientorange, very thick, fill=quotientorange!20] coordinates {
  (-0.4,0) (0,0) (1,1) (2,0) (4.4,0)
};
\node[font=\footnotesize, quotientorange, anchor=south west] at (axis cs:1.1,0.6) {corners at $0,1,2$};
\end{axis}

\begin{axis}[
  name=p3, at={($(p1.south)-(0,1.9cm)$)}, anchor=north,
  width=7.2cm, height=4.2cm,
  xlabel={$x$}, ylabel={density},
  xmin=-0.5, xmax=4.5, ymin=0, ymax=1.15,
  title style={font=\small},
  title={Step 3: Sum of $3$ ($C^1$, quadratic spline)},
]
\addplot[structurepurple, very thick, fill=structurepurple!20] coordinates {
  (0,0)(0.2,0.02)(0.4,0.08)(0.6,0.18)(0.8,0.32)(1.0,0.50)(1.2,0.66)
  (1.4,0.74)(1.5,0.75)(1.6,0.74)(1.8,0.66)(2.0,0.50)(2.2,0.32)
  (2.4,0.18)(2.6,0.08)(2.8,0.02)(3.0,0)
};
\node[font=\footnotesize, structurepurple, anchor=south] at (axis cs:1.5,0.78) {rounded crest};
\end{axis}

\begin{axis}[
  name=p4, at={($(p2.south)-(0,1.9cm)$)}, anchor=north,
  width=7.2cm, height=4.2cm,
  xlabel={$z$ (standardized)}, ylabel={density},
  xmin=-3.2, xmax=3.2, ymin=0, ymax=0.58,
  ytick={0,0.2,0.4},
  title style={font=\small},
  title={Step 4: Standardized $n=4$ vs.\ Gaussian ($C^\infty$)},
  legend style={font=\scriptsize, at={(0.96,0.96)}, anchor=north east},
]
\addplot[densityred, very thick, domain=-3.1:3.1, samples=80] {0.3989*exp(-x^2/2)};
\addlegendentry{$\Normal(0,1)$ target}
\addplot[sourceblue, thick, dashed, domain=-3.1:3.1, samples=60] {0.395*exp(-x^2/2.04)};
\addlegendentry{Standardized $n=4$}
\end{axis}
\end{tikzpicture}
\caption{The iterative convolution smoothing progression. Adding independent uniform random variables $U_1 + \dots + U_n$ convolved step by step: $n=1$ is a piecewise constant box ($C^{-1}$ with jump steps); $n=2$ integrates to a piecewise linear triangle ($C^0$); $n=3$ integrates to a continuously differentiable quadratic spline ($C^1$); by $n=4$, the standardized density is visually indistinguishable from the Gaussian bell curve ($C^\infty$). Every convolution smooths away sharp features.}
\label{ch21:fig-convolution-smoothing}
\end{figure}

\begin{historybox}
The bell curve was discovered in stages by mathematicians trying to solve very practical
calculation bottlenecks. In 1733, French-born mathematician Abraham de Moivre was
consulting for gamblers and trying to calculate exact binomial probabilities for long series
of coin flips. Computing factorials like $\binom{1000}{500}$ was humanly impossible by hand.
Armed with what we now call Stirling's approximation (an asymptotic formula replacing large
factorials with powers and square roots, which de Moivre co-developed), he discovered that
the middle binomial terms hugged a smooth curve proportional to $e^{-x^2/2}$. This was the
very first central limit theorem---the \emph{de Moivre--Laplace theorem}---though nobody named
it that yet.

Nearly eighty years later, in 1810, Pierre-Simon Laplace generalized de Moivre's coin-flip
result to sums of continuous measurement errors in astronomy and census surveys. Laplace
invented characteristic functions (transforming probability distributions into complex
exponentials to turn sums into products). Yet Laplace's derivation sparked a century of
controversy. His handling of limits and complex integrals was brilliant but heuristic,
interchanging limits freely without specifying when the theorem held. For decades,
mathematicians debated whether the bell curve was a universal law of nature or an artifact
of symmetric, convenient examples. Augustin-Louis Cauchy and Sim\'eon Denis Poisson argued
over whether heavy tails could break it, but lacked the analytical machinery to prove it.

Rigorous proof arrived from the St.~Petersburg school in Russia. Pafnuty Chebyshev introduced
the method of moments (matching all integer powers $\E[X^k]$ to the normal's moments), his
student Andrei Markov refined it, and in 1901 Aleksandr Lyapunov delivered the first fully
rigorous proof using characteristic functions and precise epsilon-delta bounds. In 1922,
German mathematician Jarl Waldemar Lindeberg identified the clean condition guaranteeing that
no single term dominates the sum, which William Feller later proved was essentially necessary.
The name ``central limit theorem'' was coined by George P\'olya in 1920; he meant central in the
sense of \emph{centrally important to probability theory}, not central as in the middle.

The theorem underpins modern statistical inference, Gauss's method of least squares,
experimental error analysis, and statistical mechanics (Maxwell and Boltzmann's velocity
distributions for gas molecules). Today, active frontiers remain: closing the gap on the
exact optimal constant in the Berry--Esseen finite-sample bound (the best proven upper bound
is $0.4748$, while the theoretical lower bound is $\approx 0.4097$), establishing
high-dimensional central limit theorems where the number of features $d$ grows alongside
the sample size $n$ in machine learning, and proving universality for eigenvalues of large
random matrices.
\end{historybox}

\section{Standardization: stopping the curve from running away}
\label{ch21:standardization}

Before we can state the theorem we have to deal with a bookkeeping problem. If
$X_1, X_2, \dots$ are independent and identically distributed (i.i.d.---each draw is
sampled from the exact same distribution, completely unaffected by the other draws) with
mean $\mu$ and variance $\sigma^2$, then the raw sum
\[
  S_n = X_1 + X_2 + \cdots + X_n
\]
has mean $n\mu$ and variance $n\sigma^2$. Both of those grow without bound. The center
of $S_n$ marches off to infinity at rate $n$, and its spread widens at rate $\sqrt{n}$.
So $S_n$ itself has no limiting distribution at all -- it just drifts away and fattens.
We have to chase it with a coordinate change.

Two natural normalizations clarify what scaling is needed.

\textbf{Fix the center: the sample mean.} Divide by $n$:
\[
  \bar X_n = \frac{S_n}{n}, \qquad \E[\bar X_n] = \mu, \qquad \Var(\bar X_n) = \frac{\sigma^2}{n}.
\]
Now the center is pinned at $\mu$, but the variance shrinks to zero. By the Law of Large
Numbers $\bar X_n \to \mu$. The mean stops running away, but it collapses to a single
point spike -- it loses all its shape. We have over-corrected.

\textbf{Fix both: standardize.} Subtract the mean and divide by the standard deviation.
Define
\[
  Z_n = \frac{S_n - n\mu}{\sigma\sqrt{n}}
      = \frac{\bar X_n - \mu}{\sigma/\sqrt{n}}.
\]
By construction $\E[Z_n] = 0$ and $\Var(Z_n) = 1$ for every $n$. We have removed the
drift and frozen the spread. Now $Z_n$ is neither running away nor collapsing; it is
sitting still, mean $0$ and variance $1$, ready to converge to something with a fixed
shape. That something is the standard normal.

\begin{recoverybox}
\textbf{Mental check (pause and think):} Suppose you quadruple your sample size from $n$ to $4n$.
What happens to the spread of the raw sum $S_n$? What happens to the spread of the sample mean $\bar X_n$?

\emph{Answer:} The raw sum's standard deviation $\sigma\sqrt{n}$ doubles (it increases by $\sqrt{4} = 2$), because adding more terms creates more total variability. But the sample mean's standard deviation $\sigma/\sqrt{n}$ is cut in half (it shrinks by $1/\sqrt{4} = 1/2$), because dividing the sum by $4n$ crushes the fluctuations faster than the sum generates them.
\end{recoverybox}

\begin{focusbox}
The key factor is $\sqrt{n}$ in the denominator. The sum's fluctuations around its
mean are of size $\sigma\sqrt{n}$ -- not $n$, not $\sqrt[3]{n}$, but exactly $\sqrt{n}$.
This scaling governs estimation throughout statistics: to halve the
uncertainty about an average, one must \emph{quadruple} the sample size.
\end{focusbox}

\begin{whybox}
Why $\sqrt{n}$ and not something else? Variance adds for independent quantities:
$\Var(S_n) = \Var(X_1) + \cdots + \Var(X_n) = n\sigma^2$. Standard deviation is the
square root of variance, so $\SD(S_n) = \sigma\sqrt{n}$. The independence is essential.
If the $X_i$ were perfectly correlated (all equal), the sum would just be $nX_1$,
with standard deviation $n\sigma$ -- linear in $n$, not square-root. The $\sqrt{n}$ scaling
is a direct consequence of independence: larger samples yield more precise estimates
precisely because uncorrelated fluctuations cancel each other out.
\end{whybox}

\section{The statement}
\label{ch21:statement}

\begin{theorem}[Central Limit Theorem, i.i.d. version]
\label{ch21:thm-clt}
Let $X_1, X_2, \dots$ be i.i.d.\ random variables with finite mean $\mu = \E[X_1]$ and
finite, nonzero variance $\sigma^2 = \Var(X_1)$. Let $S_n = \sum_{i=1}^n X_i$. Then the
standardized sum converges in distribution to a standard normal:
\[
  Z_n = \frac{S_n - n\mu}{\sigma\sqrt{n}} \;\dto\; \Normal(0,1).
\]
Equivalently, for every real number $z$,
\[
  \P\!\left( \frac{S_n - n\mu}{\sigma\sqrt{n}} \le z \right) \longrightarrow \Phi(z)
  \quad \text{as } n \to \infty,
\]
where $\Phi(z) = \int_{-\infty}^{z} \frac{1}{\sqrt{2\pi}} e^{-t^2/2}\,\dd t$ is the
standard normal cumulative distribution function.
\end{theorem}

Three details in this statement are doing the real work.

\begin{itemize}
\item \textbf{The hypotheses are minimal.} We need i.i.d., a finite mean, and a finite
nonzero variance. We need \emph{nothing} about the shape of the $X_i$. They can be
discrete or continuous, symmetric or skewed, bounded or unbounded. As long as the
variance is finite, the bell appears.
\item \textbf{The conclusion is universal.} The limit is \emph{always} the same standard
normal $\Normal(0,1)$, no matter what the $X_i$ were. The original distribution's
identity survives only through the two numbers $\mu$ and $\sigma^2$ used to standardize.
\item \textbf{Convergence is in distribution.} This means the CDFs converge,
$\P(Z_n \le z) \to \Phi(z)$, pointwise at every $z$ (here, every $z$, because $\Phi$ is
continuous). It does \emph{not} say $Z_n$ converges to any particular random number. It
is a statement about probabilities and shapes, not about realized values.
\end{itemize}

\begin{trapbox}
The finite-variance hypothesis is not a technicality you can wave away. If the $X_i$ have
infinite variance -- like a Cauchy distribution, whose tails are so heavy that
$\E[X^2] = \infty$ and even $\E[X]$ fails to exist -- the CLT is \emph{false}. Instead of approaching a normal curve, the
average of $n$ independent standard Cauchy variables has the exact same
Cauchy distribution for every $n$. Averaging does not help at all. Heavy tails
break the theorem because a single extreme observation can dominate the whole sum,
preventing independent fluctuations from averaging out. There is a beautiful generalization (the
\emph{stable laws}, with $\alpha$-stable limits) for these cases, but the friendly normal
limit requires finite variance.
\end{trapbox}

In the language of the sample mean, dividing top and bottom by $n$ gives the form you
will actually use in statistics:
\[
  \frac{\bar X_n - \mu}{\sigma/\sqrt{n}} \;\dto\; \Normal(0,1),
  \qquad \text{informally} \qquad
  \bar X_n \;\approx\; \Normal\!\left(\mu, \frac{\sigma^2}{n}\right).
\]
The sample mean of a large sample is approximately normal, centered at the true mean,
with standard deviation $\sigma/\sqrt{n}$ -- the famous \emph{standard error}.

\section{Watching it happen}
\label{ch21:simulation}

The fastest way to believe the CLT is to simulate it. Take a deliberately ugly starting
distribution and watch its standardized sums smooth into a bell. We will use a very
skewed one: the exponential distribution, which has a hard wall at $0$ and a long right
tail. It looks nothing like a bell.

Figure~\ref{ch21:fig-converge} overlays the density of the standardized sum
$Z_n = (S_n - n\mu)/(\sigma\sqrt n)$ for an exponential starting variable, at $n=1$,
$n=3$, and $n=30$, against the target standard-normal curve. At $n=1$ it is the raw
exponential shape (skewed, capped on the left). By $n=3$ the skew is visibly draining
away. By $n=30$ it is nearly indistinguishable from the bell.

\begin{figure}[H]
\centering
\begin{tikzpicture}
\begin{axis}[
  width=13cm, height=6cm,
  xlabel={standardized value $z$}, ylabel={density},
  xmin=-3.5, xmax=4.2, ymin=0, ymax=0.55,
  title style={font=\small},
  title={Standardized sum of $n$ exponentials converging to $\Normal(0,1)$},
  legend style={font=\footnotesize, at={(0.98,0.97)}, anchor=north east},
  domain=-3.5:4.2, samples=180,
]
% target standard normal
\addplot[densityred, very thick] {0.3989422804*exp(-x^2/2)};
\addlegendentry{$\Normal(0,1)$ target}
% n=1: standardized Exp(1): density of Z = X-1 for x>=-1 is e^{-(z+1)}
\addplot[sourceblue, thick, domain=-1:4.2] {exp(-(x+1))};
\addlegendentry{$n=1$ (raw exponential)}
% n=3: standardized Gamma(3,1): z=(s-3)/sqrt3, s=sqrt3 z +3, jacobian sqrt3
\addplot[quotientorange, thick, domain=-1.732:4.2, samples=160]
  {1.732050808*((1.732050808*x+3)^2)*exp(-(1.732050808*x+3))/2};
\addlegendentry{$n=3$}
% n=30: exact standardized Gamma(30,1) density
\addplot[targetgreen, thick, dashed, smooth] coordinates {
  (-3.20,0.00014) (-2.90,0.00098) (-2.61,0.00457) (-2.31,0.01559)
  (-2.02,0.04117) (-1.72,0.08790) (-1.42,0.15667) (-1.13,0.23911)
  (-0.83,0.31897) (-0.54,0.37814) (-0.24,0.40390) (0.06,0.39320)
  (0.35,0.35225) (0.65,0.29280) (0.94,0.22743) (1.24,0.16608)
  (1.54,0.11464) (1.83,0.07514) (2.13,0.04696) (2.42,0.02809)
  (2.72,0.01613) (3.02,0.00892) (3.31,0.00476) (3.61,0.00246)
  (3.90,0.00123) (4.20,0.00060)
};
\addlegendentry{$n=30$ (essentially the bell)}
\end{axis}
\end{tikzpicture}
\caption{The exponential distribution is strongly right-skewed, yet its standardized
sums march steadily toward the standard normal. At $n=1$ the curve is the raw skewed
exponential; by $n=3$ most of the skew is gone; by $n=30$ it is visually the bell. The
left wall of the exponential is what takes longest to erode.}
\label{ch21:fig-converge}
\end{figure}

\begin{happeningbox}
Observe the progression as $n$ grows in
Figure~\ref{ch21:fig-converge}. The exponential distribution is bounded below: a single draw
can never be more than $1$ standard deviation below its mean (since the minimum value
is $0$ and the mean is $1$). When standardizing a single draw, nothing can fall left
of $z=-1$. But the \emph{sum} of many draws can dip below its mean by several standard
deviations, because the lower bound on each individual term has diminishing influence on the sum.
As $n$ increases, the support of the standardized sum expands to the left, matching the
normal distribution's symmetric tail. The skewness decreases more slowly than the variance
stabilizes, which is quantified precisely by the Berry--Esseen bound in
Section~\ref{ch21:berry-esseen}.
\end{happeningbox}

\begin{labbox}
Try this yourself with a few lines of code (pseudocode; any language works):
\begin{enumerate}
\item Pick a wild distribution: maybe \texttt{X = (a uniform)\^{}4}, or a 90/10 mixture
of a spike at $0$ and a uniform on $[0,10]$. Compute its true $\mu$ and $\sigma$
(estimate them from a million draws if you must).
\item For each of $n \in \set{1, 2, 5, 30, 100}$: simulate $S_n$ a hundred thousand
times, standardize each replicate to $z = (S_n - n\mu)/(\sigma\sqrt n)$, and histogram
the $z$ values.
\item Overlay the standard normal density. Watch the histogram crawl toward the bell as
$n$ grows.
\end{enumerate}
Then deliberately break it: use a Cauchy or a Pareto with index $\le 2$ (infinite
variance) and watch the histogram \emph{refuse} to converge. Seeing the failure case is
as instructive as seeing the success.
\end{labbox}

\section{Why it is true: the moment generating function sketch}
\label{ch21:why}

The simulations make the CLT believable. Seeing \emph{why} it must be true takes a
short argument, and the cleanest one uses the moment generating function (MGF). It is not fully
rigorous -- the airtight proof uses characteristic functions, which are MGFs with an
imaginary argument and which always exist -- but the MGF version shows you the whole
mechanism with ordinary calculus, and the upgrade to characteristic functions is purely
technical.

Recall the MGF of a random variable $Y$ is $M_Y(t) = \E[e^{tY}]$, when it exists in a
neighborhood of $t=0$. Two facts make MGFs the perfect tool here.

\begin{itemize}
\item \textbf{MGFs turn sums into products.} If $Y_1, \dots, Y_n$ are independent, then
$M_{Y_1 + \cdots + Y_n}(t) = \prod_i M_{Y_i}(t)$, because the expectation of a product
of independent things factors. Summing random variables, which is hard, becomes
multiplying their MGFs, which is easy.
\item \textbf{MGFs determine distributions, and convergence of MGFs implies
convergence in distribution.} This is the \emph{continuity theorem}: if
$M_{Z_n}(t) \to M_Z(t)$ for all $t$ near $0$, and $M_Z$ is the MGF of some distribution,
then $Z_n \dto Z$. So if we can show the standardized sum's MGF converges to the normal's
MGF, we are done.
\end{itemize}

The standard normal's MGF is the target we are aiming at:
\[
  M_{\Normal(0,1)}(t) = \E[e^{tZ}] = e^{t^2/2}.
\]
Our whole job is to show $M_{Z_n}(t) \to e^{t^2/2}$.

\textbf{Setup.} Without loss of generality center and scale the $X_i$ so that $\mu = 0$
and $\sigma = 1$ (just replace $X_i$ by $(X_i - \mu)/\sigma$; this changes nothing about
the shape of the limit). Then $S_n = X_1 + \cdots + X_n$ and
$Z_n = S_n / \sqrt{n}$. Let $M(t) = \E[e^{tX_1}]$ be the common MGF of the centered,
scaled variables. Because $\E[X_1] = 0$ and $\E[X_1^2] = 1$, the Taylor expansion of $M$
around $t=0$ begins
\[
  M(t) = \E\!\left[ 1 + tX_1 + \tfrac{t^2}{2} X_1^2 + \cdots \right]
       = 1 + t\cdot 0 + \tfrac{t^2}{2}\cdot 1 + o(t^2)
       = 1 + \tfrac{t^2}{2} + o(t^2).
\]
This is the key point: \emph{the first two derivatives of $M$ at $0$ are
the mean and the second moment}, and after centering and scaling those are $0$
and $1$. All higher moments of $X_1$ enter only through the $o(t^2)$ remainder,
which vanishes in the limit.

\textbf{Multiply, because the terms are independent.} Each $X_i / \sqrt n$ has MGF
$M(t/\sqrt n)$, so
\[
  M_{Z_n}(t) = \E\!\left[ \exp\!\Big( \tfrac{t}{\sqrt n}\textstyle\sum_i X_i \Big) \right]
             = \prod_{i=1}^n M\!\left( \tfrac{t}{\sqrt n} \right)
             = \left[ M\!\left( \tfrac{t}{\sqrt n} \right) \right]^{n}.
\]
Now substitute the Taylor expansion with the small argument $t/\sqrt n$:
\[
  M\!\left( \tfrac{t}{\sqrt n} \right)
   = 1 + \frac{1}{2}\!\left(\frac{t}{\sqrt n}\right)^{\!2} + o\!\left(\frac{1}{n}\right)
   = 1 + \frac{t^2}{2n} + o\!\left(\frac{1}{n}\right).
\]
Raise to the $n$-th power and take the limit. To see why the little $o(1/n)$ disappears,
take the natural logarithm:
\[
  \ln M_{Z_n}(t) = n \ln\!\left( 1 + \frac{t^2/2}{n} + o\!\left(\tfrac1n\right) \right)
                 = n \left( \frac{t^2/2}{n} + o\!\left(\tfrac1n\right) \right)
                 = \frac{t^2}{2} + o(1).
\]
As $n \to \infty$, the $o(1)$ term vanishes, leaving exactly $t^2/2$. Exponentiating back gives
the classic $\left(1 + \frac{c}{n}\right)^n \to e^{c}$ limit:
\[
  M_{Z_n}(t) = \left[ 1 + \frac{t^2/2}{n} + o\!\left(\tfrac1n\right) \right]^{n}
             \;\longrightarrow\; e^{t^2/2}.
\]
That is exactly the standard normal MGF. By the continuity theorem,
$Z_n \dto \Normal(0,1)$. $\qed$

\begin{whybox}
This Taylor expansion explains the \emph{universality} of the theorem. The
only features of $X_1$ that survived into the limit were the coefficients of $t^0$,
$t^1$, and $t^2$ in the Taylor expansion of $M$ -- that is, $1$, the mean, and the
second moment. We forced the mean to $0$ and the variance to $1$ by standardizing. So
every legal starting distribution produces the identical leading expansion
$1 + t^2/(2n) + \cdots$. The higher moments (skewness lives in the $t^3$ term, kurtosis
in the $t^4$ term) all entered with extra powers of $1/\sqrt n$ and vanished in
the limit. The algebraic attractor $\big(1 + \tfrac{c}{n}\big)^n \to e^c$ isolates
precisely the first two moments.
\end{whybox}

\begin{formalbox}
\textbf{The rigorous characteristic function proof.}
The moment generating function $M(t) = \E[e^{tX}]$ may fail to exist if the tails of $X$
decay slower than exponential (for instance, a Student-$t$ distribution with $5$ degrees
of freedom has finite variance but infinite MGF for all $t \ne 0$). The rigorous,
universal proof replaces $M(t)$ with the \emph{characteristic function}
$\varphi(t) = \E[e^{itX}]$, which is guaranteed to exist for all $t \in \R$ because
$\abs{e^{itX}} = 1$.

Let $X_1, \dots, X_n$ be i.i.d.\ with $\E[X_1] = 0$ and $\Var(X_1) = 1$. Taylor's
theorem applied to the complex exponential $u \mapsto e^{iu}$ gives
\[
  e^{iu} = 1 + iu - \frac{u^2}{2} + R(u), \qquad \text{where } \abs{R(u)} \le \min\left(u^2, \frac{\abs{u}^3}{6}\right).
\]
Substituting $u = t X_1$ and taking expectations yields
\[
  \varphi(t) = \E[e^{itX_1}] = 1 + i t \E[X_1] - \frac{t^2}{2}\E[X_1^2] + \E[R(tX_1)]
             = 1 - \frac{t^2}{2} + o(t^2) \quad \text{as } t \to 0,
\]
since $\frac{1}{t^2}\E[\abs{R(tX_1)}] \le \E[\min(X_1^2, \tfrac{|t|}{6}|X_1|^3)] \to 0$ as $t \to 0$ by the Dominated Convergence Theorem (the integrand vanishes pointwise and is bounded by the integrable ceiling $X_1^2$, where $\E[X_1^2] = 1 < \infty$).
Now consider the standardized sum $Z_n = \frac{1}{\sqrt n} \sum_{k=1}^n X_k$. By independence,
\[
  \varphi_{Z_n}(t) = \E\!\left[ \exp\!\Big( i t \frac{1}{\sqrt n}\sum_{k=1}^n X_k \Big) \right]
                   = \left[ \varphi\!\left(\frac{t}{\sqrt n}\right) \right]^n
                   = \left[ 1 - \frac{t^2}{2n} + o\!\left(\frac{1}{n}\right) \right]^n.
\]
Taking the principal branch of the complex logarithm,
\[
  n \log\left( 1 - \frac{t^2}{2n} + o\!\left(\frac{1}{n}\right) \right)
  = n \left( -\frac{t^2}{2n} + o\!\left(\frac{1}{n}\right) \right)
  = -\frac{t^2}{2} + o(1).
\]
Exponentiating both sides:
\[
  \lim_{n\to\infty} \varphi_{Z_n}(t) = e^{-t^2/2} \quad \text{pointwise for every } t \in \R.
\]
The function $e^{-t^2/2}$ is continuous at $t=0$ and is precisely the characteristic function of
a standard normal variable $\Normal(0,1)$. By L\'evy's Continuity Theorem, pointwise convergence
of characteristic functions to a continuous limit implies convergence in distribution:
$Z_n \dto \Normal(0,1)$.
\end{formalbox}

\begin{connectionbox}
The limit $\big(1 + \frac{c}{n}\big)^n \to e^c$ is the same one that defines the number
$e$ and that drives continuous compound interest in calculus -- see the companion
\emph{Calculus From the Ground Up} for the careful $\eps$--$N$ version. It is striking
that the Central Limit Theorem, the law of measurement error, and the exponential
function all rest on the same little limit. The function $e^{-t^2/2}$ is the unique characteristic
function fixed by ``mean $0$, variance $1$, and maximal entropy,'' which is another way of saying the
normal is the unique distribution remembering only two moments.
\end{connectionbox}

\section{The de Moivre--Laplace special case: coins}
\label{ch21:demoivre}

The oldest and most concrete CLT is for coin flips. Let $X_i$ be i.i.d.\ $\Bern(p)$, so
$S_n \sim \Bin(n, p)$ counts successes in $n$ trials. Here $\mu = p$ and
$\sigma^2 = p(1-p)$, so the CLT specializes to
\[
  \frac{S_n - np}{\sqrt{n p (1-p)}} \;\dto\; \Normal(0,1).
\]
This is the de Moivre--Laplace theorem, and it is enormously useful because binomial
probabilities are tedious to compute exactly for large $n$ but trivial to approximate
with a normal.

\begin{example}[A poll]
\label{ch21:ex-poll}
You poll $n = 1000$ voters; suppose the true support for a candidate is $p = 0.52$. What
is the probability your sample shows a \emph{minority} -- that is, $\bar X_n \le 0.50$?

Here $\mu = 0.52$, $\sigma = \sqrt{0.52 \cdot 0.48} = \sqrt{0.2496} \approx 0.4996$, and
the standard error of the mean is $\sigma/\sqrt n = 0.4996/\sqrt{1000} \approx 0.0158$.
Standardize the threshold:
\[
  z = \frac{0.50 - 0.52}{0.0158} \approx -1.266.
\]
So $\P(\bar X_n \le 0.50) \approx \Phi(-1.266) \approx 0.103$. Even with a true majority,
roughly a one-in-ten poll of $1000$ people would show the candidate behind. Want to cut
that to one in a hundred? You would need $z \approx -2.33$, which means shrinking the
standard error by a factor of $2.33/1.266 \approx 1.84$, which by the $\sqrt n$ law
requires about $1.84^2 \approx 3.4$ times the sample, i.e.\ around $3400$ voters. The
square-root tax in action.
\end{example}

\begin{figure}[H]
\centering
\begin{tikzpicture}
\begin{axis}[
  width=13cm, height=5.5cm,
  ybar, bar width=5pt,
  xlabel={number of heads}, ylabel={probability},
  xmin=12, xmax=38, ymin=0, ymax=0.13,
  scaled y ticks=false,
  yticklabel style={/pgf/number format/fixed, /pgf/number format/precision=2},
  title style={font=\small},
  title={$\Bin(50,0.5)$ bars vs.\ its normal approximation},
  legend style={font=\footnotesize, at={(0.02,0.97)}, anchor=north west},
]
% Binomial(50,0.5) pmf near the center (approx values)
\addplot+[draw=sourceblue!60, fill=histfill!70] coordinates {
 (15,0.0020)(17,0.0087)(19,0.0270)(21,0.0598)(23,0.0960)
 (25,0.1123)(27,0.0960)(29,0.0598)(31,0.0270)(33,0.0087)(35,0.0020)};
\addlegendentry{$\Bin(50,0.5)$ pmf (odd values)}
% Normal(25, 12.5): density 1/sqrt(2*pi*12.5) e^{-(x-25)^2/25}
\addplot[densityred, very thick, domain=12:38, samples=120]
  {0.112838*exp(-(x-25)^2/25)};
\addlegendentry{$\Normal(25,12.5)$}
\end{axis}
\end{tikzpicture}
\caption{The binomial $\Bin(50, 0.5)$ probability mass function (bars, odd values
shown for clarity) with its normal approximation $\Normal(25, 12.5)$ (smooth curve) laid
over it. The continuous curve threads through the tops of the discrete bars; to convert
between them correctly we need the continuity correction of
Section~\ref{ch21:continuity}.}
\label{ch21:fig-binom}
\end{figure}

\section{The continuity correction}
\label{ch21:continuity}

There is a subtlety when you approximate a \emph{discrete} distribution (like the
binomial, which lives on the integers) by a \emph{continuous} one (the normal). The
binomial puts a lump of probability exactly at each integer $k$; the normal smears
probability continuously and assigns probability $0$ to any single point. To translate
accurately, you must decide how to split the gaps between integers.

The fix is the \emph{continuity correction}: represent each integer $k$ by the unit-width
interval $[k - \tfrac12, k + \tfrac12]$ when you pass to the normal. Each integer ``owns'' the
half-unit on either side of it. Concretely:
\[
  \P(S_n = k) \approx \P\!\left(k - \tfrac12 \le \Normal \le k + \tfrac12\right),
\]
\[
  \P(S_n \le k) \approx \Phi\!\left(\frac{k + \tfrac12 - n\mu}{\sigma\sqrt n}\right),
  \qquad
  \P(S_n \ge k) \approx 1 - \Phi\!\left(\frac{k - \tfrac12 - n\mu}{\sigma\sqrt n}\right).
\]
The half-unit shifts are easy to get backwards, so reason about which integers you mean
to \emph{include}. ``At most $k$'' includes $k$, so push the boundary up to $k+\tfrac12$
to capture the entire bar. ``At least $k$'' includes $k$, so push the boundary down to $k-\tfrac12$.
Figure~\ref{ch21:fig-continuity} shows the geometry behind this shift.

\begin{figure}[H]
\centering
\begin{tikzpicture}
\begin{axis}[
  width=13cm, height=5.8cm,
  xlabel={value $x$}, ylabel={probability / density},
  xmin=6.2, xmax=13.8, ymin=0, ymax=0.29,
  xtick={7,8,9,10,11,12,13},
  title style={font=\small},
  title={Geometry of the continuity correction for $\P(S_n \le k)$ with $k=10$},
  axis lines=left,
]
\addplot[draw=none, fill=densityred!15, domain=6.2:10.5, samples=60] {0.1735 * exp(-(x-10)^2 / (2 * 2.3^2))} \closedcycle;
\addplot[densityred, very thick, domain=6.2:13.8, samples=120] {0.1735 * exp(-(x-10)^2 / (2 * 2.3^2))};
\draw[draw=sourceblue, fill=sourceblue!30] (axis cs:6.5, 0) rectangle (axis cs:7.5, 0.045);
\draw[draw=sourceblue, fill=sourceblue!30] (axis cs:7.5, 0) rectangle (axis cs:8.5, 0.090);
\draw[draw=sourceblue, fill=sourceblue!30] (axis cs:8.5, 0) rectangle (axis cs:9.5, 0.145);
\draw[draw=sourceblue, fill=sourceblue!35] (axis cs:9.5, 0) rectangle (axis cs:10.5, 0.170);
\draw[draw=gray!60, fill=gray!10] (axis cs:10.5, 0) rectangle (axis cs:11.5, 0.145);
\draw[draw=gray!60, fill=gray!10] (axis cs:11.5, 0) rectangle (axis cs:12.5, 0.090);
\draw[draw=gray!60, fill=gray!10] (axis cs:12.5, 0) rectangle (axis cs:13.5, 0.045);
\draw[densityred, very thick, dashed] (axis cs:10.5, 0) -- (axis cs:10.5, 0.185);
\draw[gray!70!black, thick, dashdotted] (axis cs:10, 0) -- (axis cs:10, 0.185);
\draw[<-, thick, densityred] (axis cs:10.5, 0.185) -- (axis cs:12.0, 0.22)
  node[above, font=\footnotesize, align=center] {Correct boundary $k+0.5=10.5$\\[2pt]Includes entire bar for $k$};
\draw[<-, thick, gray!70!black] (axis cs:10.0, 0.185) -- (axis cs:8.2, 0.22)
  node[above, font=\footnotesize, align=center] {Uncorrected boundary $k=10$\\[2pt]Cuts bar $k$ in half};
\end{axis}
\end{tikzpicture}
\caption{Why the continuity correction adds $\frac12$. Each discrete integer $k$ is represented by a bar of width $1$ centered at $k$, covering the interval $[k-\frac12, k+\frac12]$. To evaluate $\P(S_n \le k)$ using the smooth normal curve, the continuous integral must extend to the right edge of the $k$-th bar at $k+\frac12$; stopping at $k$ would lose half of that bar's probability mass.}
\label{ch21:fig-continuity}
\end{figure}

\begin{example}[Continuity correction pays off]
\label{ch21:ex-cc}
Flip a fair coin $n = 100$ times. What is $\P(S_n \le 45)$? Here $\mu = 0.5$,
$\sigma = 0.5$, so $n\mu = 50$ and $\sigma\sqrt n = 0.5 \cdot 10 = 5$.

Without the correction:
$\Phi\!\big(\frac{45 - 50}{5}\big) = \Phi(-1.0) \approx 0.1587$.

With the correction:
$\Phi\!\big(\frac{45.5 - 50}{5}\big) = \Phi(-0.9) \approx 0.1841$.

The exact binomial answer is $0.1841$ to four places. The corrected approximation nails
it; the uncorrected one is off by nearly three percentage points. The smaller your $n$
or the closer you are to the tails, the more the half-unit matters.
\end{example}

\begin{trapbox}
Two common slips. First, forgetting the correction entirely when the problem is plainly
discrete -- this is the single most frequent CLT error on exams. Second, applying it to
a genuinely continuous quantity, where it does nothing but add noise: if your $X_i$ are
continuous (waiting times, measurements), there are no integer lumps and no correction
to make. The continuity correction is for discrete-to-continuous translation only.
\end{trapbox}

\begin{recoverybox}
\textbf{Mental check (pause and think):} You roll a fair die $60$ times. Let $K$ be the number of sixes. You want to approximate $\P(K > 12)$ with the normal distribution. What boundary value should you plug into the standardized formula?

\emph{Answer:} Because $K$ counts discrete rolls, the strict inequality $K > 12$ is identical to $K \ge 13$. The included integers start at $13$, so the continuity correction moves the boundary down by half a unit to $12.5$. We evaluate the normal tail $1 - \Phi$ at $12.5$. If you had carelessly shifted $12$ up to $12.5$, you would have gotten the right number by accident---but for $\P(K < 12) = \P(K \le 11)$, the boundary moves up to $11.5$! Always translate strict inequalities to inclusive ones ($\le$ or $\ge$) before shifting by $\pm \frac12$.
\end{recoverybox}

\section{How fast? The Berry--Esseen theorem}
\label{ch21:berry-esseen}

The CLT says $\P(Z_n \le z) \to \Phi(z)$, but it is silent about \emph{how fast}. For a
nearly symmetric, light-tailed starting distribution, $n = 20$ may already give an
excellent fit; for a violently skewed one, even $n = 200$ may leave a visible error in
the tails. We would like a guarantee. Berry (1941) and Esseen (1942) provided one.

\begin{theorem}[Berry--Esseen]
\label{ch21:thm-be}
Let $X_1, \dots, X_n$ be i.i.d.\ with mean $\mu$, variance $\sigma^2 > 0$, and finite
third absolute moment $\rho = \E\big[\,\abs{X_1 - \mu}^3\,\big] < \infty$. Let $F_n$ be
the CDF of the standardized sum $Z_n$. Then there is a universal constant $C$ such that
for all $n$ and all $z$,
\[
  \sup_{z \in \R} \big| F_n(z) - \Phi(z) \big| \;\le\; \frac{C\,\rho}{\sigma^3 \sqrt n}.
\]
The best known constant is $C < 0.4748$.
\end{theorem}

Two things to read off this beautiful inequality.

\textbf{The rate is $1/\sqrt n$.} The worst-case error between the true CDF and the
normal CDF shrinks like $1/\sqrt n$ -- the same $\sqrt n$ that governs everything else
in this chapter. Quadruple the sample, halve the approximation error.

\textbf{The constant scales with skewness.} The factor $\rho/\sigma^3$ is essentially
the absolute skewness of the starting distribution. A symmetric, concentrated
distribution has small $\rho/\sigma^3$ and converges fast; a lopsided, heavy-shouldered
one has large $\rho/\sigma^3$ and converges slowly. This is the precise version of the
intuition from the exponential simulation: \emph{skewness is what takes longest to wash
out}.

\begin{example}[Berry--Esseen for a coin]
\label{ch21:ex-be}
For a fair coin, $X_1 \in \set{0,1}$ with $\mu = \tfrac12$, $\sigma = \tfrac12$, and
$\rho = \E\,\abs{X_1 - \tfrac12}^3 = (\tfrac12)^3 = \tfrac18$. So
$\rho/\sigma^3 = (1/8)/(1/8) = 1$, and the bound reads
$\sup_z |F_n - \Phi| \le 0.4748/\sqrt n$. At $n = 100$ that is about $0.047$ -- a
worst-case CDF error under five percentage points, and in practice (away from the
extreme tails) far better, which is why the continuity-corrected answer in
Example~\ref{ch21:ex-cc} was essentially exact.
\end{example}

\begin{aestheticbox}
The plain CLT is an asymptotic statement: it guarantees convergence as $n \to \infty$,
but does not specify the approximation error for a given finite sample size such as $n = 30$.
Berry--Esseen provides an explicit, computable error bound for finite $n$,
depending only on three quantities ($\sigma$, $\rho$, and $n$).
Such non-asymptotic bounds are essential in statistical practice and theoretical computer science,
where explicit finite-sample guarantees are required.
\end{aestheticbox}

\section{Why the normal, and not some other curve?}
\label{ch21:universality}

Pause on the deepest question the CLT raises: of all possible shapes, why
is \emph{this} one -- $e^{-x^2/2}$ -- the universal attractor? Three complementary
answers.

\textbf{The fixed-point answer.} Suppose the limit of standardized sums exists and call
its law $L$. If you take two independent copies of $L$ and add them (rescaling to keep
variance $1$), you must get $L$ back -- because a sum of two big sums is just a bigger
sum, and the limit cannot depend on how you grouped the terms. So $L$ is a
\emph{fixed point} of the ``add two independent copies and rescale'' operation. The
normal is exactly such a fixed point: if $A, B \sim \Normal(0,1)$ independently, then
$(A+B)/\sqrt 2 \sim \Normal(0,1)$ again. The bell is the shape that reproduces itself
under averaging, and the CLT says it is the stable one that everything flows toward.

\textbf{The maximum-entropy answer.} Among all distributions with a given mean and
variance, the normal is the one of maximum entropy -- the most spread-out distribution
consistent with those two constraints. Summing independent random variables combines
their variability while preserving only the total mean and variance, so the distribution
settles toward the maximum-entropy distribution matching those constraints: the Gaussian.

\textbf{The moment answer.} We saw it in the proof: the only moments that survive
standardization and the $n \to \infty$ limit are the mean and variance. Any limit law is
therefore pinned down by mean $0$ and variance $1$ alone, and the normal is the unique
distribution determined by exactly those.

\begin{paradoxbox}
\textbf{The disappearance of distribution-specific details.} It may seem surprising that the
limit retains so few features of the original distribution. Two starting distributions can have
very different shapes -- a coin flip, an exponential variable, and a uniform variable -- yet
their standardized sums all approach the same standard normal curve.
These differences do not vanish instantly: for finite $n$, sums of Bernoulli variables and sums of
exponentials remain distinguishable by their skewness and higher moments.
What the CLT and Berry--Esseen state is that these differences shrink at rate $1/\sqrt n$
(for skewness) and faster (for kurtosis), while the Gaussian profile remains of order $1$.
In the standardized limit, the contribution of higher moments becomes negligible.
This is also a practical caution: two processes whose standardized averages both look approximately normal
may still differ substantially in their extreme tails, where rare events occur.
\end{paradoxbox}

\section{Applications: the normal everywhere}
\label{ch21:applications}

\begin{applicationbox}
\textbf{Why measurement error is normal.} A single physical measurement is corrupted by
a great many tiny, independent perturbations -- thermal jitter, air currents, electronic
noise, rounding, vibration. Each contributes a small random nudge; the total error is
their sum. By the CLT, that sum is approximately normal, with mean $0$ if the
perturbations are unbiased. This is why Gauss introduced the normal distribution
precisely to model astronomical measurement error (it is still sometimes called the
Gaussian), and why ``assume the noise is Gaussian'' is the default in physics,
engineering, and signal processing. When noise arises from many small independent disturbances, the CLT justifies that Gaussian model.
\end{applicationbox}

\begin{applicationbox}
\textbf{Confidence intervals and the standard error.} The single most-used consequence
of the CLT in all of data analysis. Given a sample of size $n$ with sample mean
$\bar X_n$ and sample standard deviation $s$, the CLT (plus a little extra work to
replace $\sigma$ by $s$) gives the approximate $95\%$ confidence interval
\[
  \bar X_n \pm 1.96 \cdot \frac{s}{\sqrt n}.
\]
The $1.96$ is the $z$-value cutting off $2.5\%$ in each normal tail; the $s/\sqrt n$ is
the standard error. Every poll's ``margin of error,'' every A/B test's significance
calculation, and every reported ``mean $\pm$ standard error'' in a scientific paper is
this formula, powered by the CLT. We develop it fully in the estimation and
hypothesis-testing chapters; here just notice that it requires \emph{nothing} about the
shape of the underlying data, only finite variance and a decent sample size.
\end{applicationbox}

\begin{applicationbox}
\textbf{A/B testing.} You show version A of a webpage to $n_A$ users and version B to
$n_B$, and measure conversion rates $\hat p_A$ and $\hat p_B$. Each rate is an average
of Bernoulli outcomes, so by de Moivre--Laplace each is approximately normal. Their
difference is then approximately normal too, with standard error
$\sqrt{\hat p_A(1-\hat p_A)/n_A + \hat p_B(1-\hat p_B)/n_B}$. Dividing the observed
difference by this standard error gives a $z$-statistic you can compare to the normal
curve to decide whether B genuinely beats A or just got lucky. The entire online
experimentation industry runs on the CLT applied to coin flips.
\end{applicationbox}

\begin{applicationbox}
\textbf{Monte Carlo simulation.} To estimate an integral or an expectation
$\theta = \E[g(X)]$ that is too hard to compute analytically, draw $n$ independent
samples and average: $\hat\theta_n = \frac1n \sum_i g(X_i)$. The Law of Large Numbers
says $\hat\theta_n \to \theta$; the CLT tells you the \emph{error}: $\hat\theta_n$ is
approximately $\Normal(\theta, \tau^2/n)$, so the Monte Carlo error shrinks like
$\tau/\sqrt n$. This $1/\sqrt n$ rate is independent of the dimension of $X$ -- which is
exactly why Monte Carlo beats grid-based numerical integration in high dimensions, and
why it is widely used in computational physics, Bayesian statistics, and quantitative
finance. See the Monte Carlo chapter for variance-reduction methods that shrink $\tau$
(the $1/\sqrt n$ rate remains, but the constant factor is reduced).
\end{applicationbox}

\begin{connectionbox}
\textbf{Where the CLT sits among the limit theorems.} The Law of Large Numbers (previous
chapter) and the Central Limit Theorem are a matched pair, and the difference is the
scaling. The LLN looks at $\bar X_n - \mu$ on its natural scale and finds it goes to
$0$: the average converges to the truth. The CLT zooms in by the factor $\sqrt n$ and
asks what the residual fluctuation $\sqrt n\,(\bar X_n - \mu)$ looks like under the
microscope -- and finds a normal. So the LLN is the \emph{first-order} fact (the average
is right) and the CLT is the \emph{second-order} fact (here is the shape of how wrong it
is). The modes-of-convergence chapter makes ``converges in distribution'' precise and
relates it to the stronger convergences the LLN uses.
\end{connectionbox}

\section{A worked example, end to end}
\label{ch21:worked}

\begin{example}[Total weight on an elevator]
\label{ch21:ex-elevator}
An elevator is rated for $1600$ kg. Suppose adult passenger weights are i.i.d.\ with
mean $\mu = 75$ kg and standard deviation $\sigma = 15$ kg (the distribution of weights
is skewed, but we never need its shape). If $20$ people board, what is the probability
the total weight exceeds the limit?

The total $S_{20}$ has mean $20 \cdot 75 = 1500$ kg and standard deviation
$\sigma\sqrt{n} = 15\sqrt{20} \approx 67.08$ kg. We want $\P(S_{20} > 1600)$.
Standardize:
\[
  z = \frac{1600 - 1500}{67.08} \approx 1.49.
\]
Weights are continuous, so no continuity correction is needed. Then
\[
  \P(S_{20} > 1600) \approx 1 - \Phi(1.49) \approx 1 - 0.9319 = 0.0681 \approx 0.068.
\]
About a $6.8\%$ chance of overload. Notice the practical utility of the CLT here: we knew only
the mean and standard deviation of an individual weight -- not the distribution's shape
at all -- and yet we produced a usable probability for the total. (How trustworthy is the $6.8\%$? With $n=20$ and a
moderately skewed weight distribution, Berry--Esseen warns the answer could be off by a
few percentage points in the tail; for a safety-critical decision you would want a
larger margin or the real weight distribution.)
\end{example}

Figure~\ref{ch21:fig-tail} shows what the calculation looks like geometrically: the
total weight is approximately normal, and the overload probability is the area in the
right tail beyond the standardized threshold $z \approx 1.49$.

\begin{figure}[H]
\centering
\begin{tikzpicture}
\begin{axis}[
  width=13cm, height=5cm,
  xlabel={standardized total weight $z$}, ylabel={density},
  xmin=-3.6, xmax=3.6, ymin=0, ymax=0.45,
  domain=-3.6:3.6, samples=160,
  ytick=\empty,
  axis lines=left,
]
% shaded right tail beyond z=1.49
\addplot[densityred, very thick, name path=bell] {0.3989422804*exp(-x^2/2)};
\addplot[draw=none, name path=axis] {0};
\addplot[densityred!25] fill between[of=bell and axis, soft clip={domain=1.49:3.6}];
\draw[dashed, gray] (axis cs:1.49,0) -- (axis cs:1.49,0.13);
\draw[<-, thick, densityred] (axis cs:1.85, 0.025) -- (axis cs:2.3, 0.08)
  node[above, font=\footnotesize, text=black] {tail area $\approx 0.068$};
\node[anchor=south, font=\footnotesize] at (axis cs:1.49, 0.14)
  {$z\approx 1.49$};
\end{axis}
\end{tikzpicture}
\caption{The elevator overload probability as a normal tail area. The total weight of
20 passengers is approximately $\Normal(1500, 67.08^2)$; the $1600$ kg limit sits
$z \approx 1.49$ standard deviations above the mean, and the shaded right tail has area
$1 - \Phi(1.49) \approx 0.068$.}
\label{ch21:fig-tail}
\end{figure}

\section{Summary}
\label{ch21:summary}

\begin{focusbox}
\textbf{Core takeaway.} For independent random variables with finite variance,
subtracting the mean $n\mu$ and dividing by $\sigma\sqrt n$ yields the standard normal limit:
\[
  \frac{S_n - n\mu}{\sigma\sqrt n} \dto \Normal(0,1).
\]
\begin{itemize}
\item \textbf{Standardization} keeps the curve centered ($-n\mu$) with constant variance
($\div\,\sigma\sqrt n$). The factor $\sqrt n$ sets the scale of statistical fluctuation.
\item \textbf{Why it works:} MGFs turn sums into products; Taylor expansion leaves only
the mean and variance; $(1 + \frac{t^2/2}{n})^n \to e^{t^2/2}$, the normal MGF. Higher
moments vanish in the limit -- hence universality.
\item \textbf{Continuity correction:} for discrete sums, replace integer $k$ by
$[k \pm \tfrac12]$ before evaluating the normal approximation.
\item \textbf{Berry--Esseen:} the CDF error is at most $0.4748\,\rho/(\sigma^3\sqrt n)$
-- rate $1/\sqrt n$, scaled by skewness.
\item \textbf{Caveat:} finite variance is required. Infinite-variance (Cauchy,
heavy-tail) distributions have non-normal stable limits.
\end{itemize}
\end{focusbox}

\subsection*{Looking ahead: what kind of convergence is this?}
The Central Limit Theorem asserts that $Z_n \dto \Normal(0,1)$. But pause on what that
arrow $\dto$ actually means. We proved that the cumulative distribution functions
$\P(Z_n \le z)$ converge pointwise to the normal CDF $\Phi(z)$. We did \emph{not} prove
that the random variable $Z_n$ gets close to any fixed normal variable $Z$ on our sample
space. In fact, $Z_n$ keeps fluctuating with every new draw, never settling down to any
single numerical value.

This distinction is the gateway to Chapter~22 (\emph{Modes of Convergence}). In Chapter~20,
the Law of Large Numbers established that $\bar X_n$ converges to $\mu$ \emph{in probability}
and \emph{almost surely}---sample averages cluster tightly around a single number.
Here in Chapter~21, the CLT establishes convergence \emph{in distribution}---the
distribution functions align, rather than the random variables themselves. In the next chapter,
we examine the complete hierarchy connecting these modes of convergence, determine when one mode
implies another, and develop core tools (the continuous mapping theorem, Slutsky's theorem,
and the delta method) that allow us to propagate normal limits through nonlinear transformations.

\section*{Exercises}
\addcontentsline{toc}{section}{Exercises}

\begin{exercise}[Mental check: absolute error vs.\ relative frequency]
\label{ch21:ex-mental-error}
A student claims: ``By the Central Limit Theorem and Law of Large Numbers, if I flip a fair
coin $n = 10{,}000$ times, the exact number of heads will almost certainly be within $5$
of $5{,}000$.'' (a) Explain the fallacy: what is the typical size of the absolute deviation
$|S_n - n/2|$ for $n = 10{,}000$? (b) Why does the \emph{fraction} of heads get closer to
$1/2$ while the \emph{count} of excess heads spreads out?
\end{exercise}

\begin{exercise}[Mental check: sample size scaling]
\label{ch21:ex-mental-scale}
You are designing an experiment whose current margin of error is $\pm 6$ units. (a) By what
factor must you multiply the sample size $n$ to shrink the margin of error to $\pm 2$ units?
(b) If the original sample had $n = 100$, what new sample size is required?
\end{exercise}

\begin{exercise}[Standardize a sum]
\label{ch21:ex-standardize}
Let $X_1, \dots, X_{40}$ be i.i.d.\ with mean $\mu = 3$ and variance $\sigma^2 = 4$. Let
$S = \sum_i X_i$. (a) Find $\E[S]$ and $\SD(S)$. (b) Using the CLT, estimate
$\P(S > 140)$. (c) Estimate $\P(115 \le S \le 130)$.
\end{exercise}

\begin{exercise}[Sample size and the square-root law]
\label{ch21:ex-sqrtn}
A quantity has standard deviation $\sigma = 8$. You want to estimate its mean with a
$95\%$ confidence interval of total width at most $1$ (i.e.\ $\pm 0.5$). Using the
interval $\bar X_n \pm 1.96\,\sigma/\sqrt n$, find the smallest $n$. Then find the $n$
needed for width $0.5$. By what factor did $n$ grow, and why?
\end{exercise}

\begin{exercise}[Continuity correction]
\label{ch21:ex-cc-prob}
A fair die is rolled $180$ times. Let $K$ be the number of sixes. (a) Find $\E[K]$ and
$\SD(K)$. (b) Use the normal approximation \emph{with} the continuity correction to
estimate $\P(K \ge 36)$. (c) Redo (b) \emph{without} the correction and report the
difference.
\end{exercise}

\begin{exercise}[A poll]
\label{ch21:ex-poll2}
In a population, $40\%$ favor a proposal. You sample $n = 600$ people. (a) What is the
approximate distribution of the sample proportion $\hat p$? (b) Estimate
$\P(\hat p \ge 0.43)$. (c) How large would $n$ need to be so that the standard error of
$\hat p$ is at most $0.01$?
\end{exercise}

\begin{exercise}[Berry--Esseen by hand]
\label{ch21:ex-be-hand}
Let $X_i$ be i.i.d.\ $\Bern(0.1)$. (a) Compute $\mu$, $\sigma^2$, and
$\rho = \E\,\abs{X_1 - \mu}^3$. (b) Compute the Berry--Esseen bound
$0.4748\,\rho/(\sigma^3\sqrt n)$ for $n = 100$ and $n = 10000$. (c) Compare to the bound
for $\Bern(0.5)$ from Example~\ref{ch21:ex-be}. Which converges to normal faster, and
does that match the intuition about skewness?
\end{exercise}

\begin{exercise}[Where the CLT fails]
\label{ch21:ex-cauchy}
The average of $n$ independent standard Cauchy random variables is again standard Cauchy,
for every $n$. (a) Explain why this does \emph{not} contradict the CLT. (b) The Cauchy
density is $f(x) = \frac{1}{\pi(1+x^2)}$. Show that $\E[X]$ does not exist by examining
$\int \abs{x} f(x)\,\dd x$. (c) In one sentence, what feature of the Cauchy breaks the
hypotheses of Theorem~\ref{ch21:thm-clt}?
\end{exercise}

\begin{exercise}[Sum of uniforms]
\label{ch21:ex-irwin}
Let $U_1, \dots, U_{12}$ be i.i.d.\ $\Unif(0,1)$. (a) Find the mean and variance of
$S = \sum_{i=1}^{12} U_i$. (b) Show that $S - 6$ has mean $0$ and variance $1$. (This is
the classic ``cheap Gaussian generator'': $\sum_{i=1}^{12} U_i - 6$ is approximately
$\Normal(0,1)$.) (c) Name one way this approximation differs from a true standard normal.
\end{exercise}

\begin{exercise}[Proof: standardization arithmetic]
\label{ch21:ex-proof-std}
Let $X_1, \dots, X_n$ be i.i.d.\ with mean $\mu$ and variance $\sigma^2$. Using only
linearity of expectation and the additivity of variance for independent variables, prove
carefully that $Z_n = (S_n - n\mu)/(\sigma\sqrt n)$ has $\E[Z_n] = 0$ and
$\Var(Z_n) = 1$ for every $n$. State exactly where independence is used.
\end{exercise}

\begin{exercise}[Proof: the MGF of a sum]
\label{ch21:ex-proof-mgf}
Let $Y_1, \dots, Y_n$ be independent with MGFs $M_{Y_i}(t)$ existing near $0$. (a) Prove
$M_{Y_1 + \cdots + Y_n}(t) = \prod_i M_{Y_i}(t)$, naming where independence enters.
(b) Specialize to the i.i.d.\ case to get $M_{S_n}(t) = M(t)^n$. (c) Using the
substitution $t \mapsto t/(\sigma\sqrt n)$, write down $M_{Z_n}(t)$ in terms of $M$ for
the centered case $\mu = 0$.
\end{exercise}

\begin{exercise}[Two-sample difference]
\label{ch21:ex-twosample}
An A/B test gives version A a conversion rate $\hat p_A = 0.118$ over $n_A = 5000$ users
and version B $\hat p_B = 0.131$ over $n_B = 5000$ users. (a) Estimate the standard error
of $\hat p_B - \hat p_A$. (b) Compute the $z$-statistic
$(\hat p_B - \hat p_A)/\text{SE}$. (c) Roughly, is the difference convincing at the
two-tailed $5\%$ level (i.e.\ is $\abs{z} > 1.96$)? Explain what the CLT is doing in
this calculation.
\end{exercise}
